\uuid{LDFg}
\titre{Primitive}
\auteur{bodin ; croizat ; sacre}
\organisation{exo7}
\licence{CC BY-NC-SA 4.0}
\datecreate{2022-03-08}
\niveau{L1}
\module{Analyse}
\chapitre{Calcul d'intégrales}
\sousChapitre{Primitives diverses}
\difficulte{3}
\theme{Primitive}
\exo7qcmid{hilisit.equations-differentielles.012}

\contenu{
\begin{question}
\qclassification{127.11}
Soit $f : I \to \Rr$ une fonction définie sur un intervalle. Soit $F$ une primitive de $f$.
$C$ désigne une constante.
Quelles sont les affirmations vraies ?
\begin{answers}
  \good{Si $f(x)=0$ sur $I$ alors $F(x)=C$.}
  \bad{Si $f(x)=x$ alors $F(x) = x^2+C$.}
  \bad{Si $f(x) \times \cos(x) = 1$ alors $F(x) = \frac{1}{\sin(x)} + C$.}
  \bad{Si $f( \ln(x) ) = 0$ alors $F(x) = e^x + C$.}
\end{answers}
\begin{explanations}
Si $f$ est la fonction nulle, alors $F$ est une fonction constante.
Si $f(x)=x$, alors $F(x) = \frac12 x^2+ C$. Les autres affirmations sont fantaisistes : lorsqu'on dérive $\frac{1}{\sin(x)} + C$ on obtient $\frac{-\cos(x)}{\sin^2(x)}$ qui n'est pas du tout l'inverse de $\cos(x)$. Et si $F(x) = e^x + C$, alors $f(x) = F'(x) = e^x$ ce qui donne $f(\ln(x)) = e^{\ln(x)} = x \neq 0$ !
\end{explanations}
\end{question}
}
